% Framing approved in the user discussion; integrated experiments remain proposed.
\section{Research capability across tools and memory}
\label{sec:framework}
We use research capability to mean selecting methods suited to a task, producing
useful decisions or analyses, and adapting later choices using available experience.
This is an evaluation target, not a property inferred from the number of tools or
successful calls. The framework separates the available methods, information retained
for the current task, and state carried across episodes.

\paragraph{A decision and learning cycle.}
Let $q_t$ describe the task, $\mathcal{D}_{\le t}$ the records available by time $t$,
and $M_t$ the persistent state. A context mechanism selects or encodes information
from the task history. A policy uses that context and a memory readout to choose a
tool and its parameters, inspect its output, and eventually produce a decision.
Once the episode ends, an eligible outcome may update $M_t$. These are conceptual
interfaces for the combined agent; the current studies implement and test subsets.

\begin{table*}[t]
\centering\small
\setlength{\tabcolsep}{5pt}
\begin{tabular}{p{0.16\textwidth}p{0.26\textwidth}p{0.25\textwidth}p{0.23\textwidth}}
\toprule
Component & Role in the research process & Capability question & Present evidence \\
\midrule
\textsc{fin-skills} & Expose guidance, method requirements and executable tools &
Which methods does the agent choose, and how does it use their outputs? &
Financial library and autonomous-use development study \\
HiSTrim & Select historical representations while protecting task constraints &
Which context can be removed while preserving decision quality? &
Reported recommendation comparisons; financial integration untested \\
Persistent memory & Read and update state using completed experiences &
Does experience change subsequent choices and improve adaptation? &
Numerical adapter and feedback interface; combined policy untested \\
Execution and feedback checks & Record actual execution and test feedback eligibility &
Which outputs and outcomes can support analysis or an update? &
Regression checks and an audit feasibility study \\
\bottomrule
\end{tabular}
\caption{Roles and evidence within the proposed financial-agent framework. The rows
describe complementary components, not an already evaluated integrated system.}
\label{tab:framework}
\end{table*}

\paragraph{Two distinct forms of memory.}
HiSTrim changes which historical representations survive within a model's attention
computation. Persistent memory changes learned state across episodes. Both can affect
future choices, but they require different interventions: a context ablation changes
what the current decision can attend to, while a memory ablation changes what previous
outcomes contribute. Neither should silently alter the information cutoff or available
tool catalog.

\paragraph{Experience as a connection to later action.}
The combined design would associate task conditions and tool outputs with a completed
outcome, then expose a memory readout at a later choice. An update counts as useful
only through its effect on subsequent tasks. A valid loss can be informative; a gain
based on unavailable information should not be treated as successful experience.
Checks test selected conditions for using feedback, while decision quality remains
an independent evaluation outcome.

Original records remain available to the checking infrastructure even if a context
selector omits them. HiSTrim requires access to a backbone's internal representations
and KV state, whereas the numerical memory adapter has a separate backend. Connecting
both to LLM tool selection requires an explicit implementation and common task suite.
